\chapter{Abdominoplasty}
\index{Abdominoplasty}

\section*{Definition and Overview}
Abdominoplasty, commonly known as a tummy tuck, is a surgical procedure that removes excess skin and fat from the abdomen and tightens the underlying abdominal wall musculature. It is performed for both cosmetic and reconstructive reasons, most often after pregnancy or major weight loss, when the skin has lost its elasticity and the rectus muscles have stretched or separated. Abdominoplasty is not a weight-loss operation and is not a substitute for diet and exercise; rather, it addresses contour, skin redundancy, and muscle laxity that cannot be corrected by lifestyle measures. The procedure is typically undertaken by specialist plastic or cosmetic surgeons, with several technical variants available to match the individual's anatomy and goals.

\section*{Epidemiology}
Abdominoplasty is among the more common cosmetic surgical procedures performed worldwide, and its popularity has increased steadily over recent decades. The majority of patients are women in their 30s to 50s, many of whom have had one or more pregnancies, but the procedure is also increasingly sought by men and by people who have undergone massive weight loss, including after bariatric surgery. Overall complication rates are relatively low when patients are carefully selected, but abdominoplasty carries greater risk than purely skin-level cosmetic procedures because of the extent of dissection and the large wounds involved.

\section*{Aetiology and Pathophysiology}
The typical candidate has sustained changes to the abdominal wall that persist despite weight stability and exercise.

\begin{itemize}
    \item \textbf{Redundant skin:} loss of skin elasticity after pregnancy or major weight loss, when the skin fails to shrink back to fit the reduced body
    \item \textbf{Diastasis recti:} widening of the linea alba, separating the paired rectus muscles and producing a midline bulge that is accentuated on straining
    \item \textbf{Persistent lower abdominal fullness} that diet and exercise cannot eliminate
    \item \textbf{Functional complaints:} skin irritation, rashes, and soreness within skin folds, plus difficulty exercising or finding well-fitting clothing
    \item \textbf{Body image concerns} that affect confidence and quality of life
\end{itemize}

\section*{Clinical Features}
A person considering abdominoplasty typically presents with a combination of the following features.

\begin{itemize}
    \item Excess, loose skin over the lower and middle abdomen
    \item Stretch marks and thinning of the skin
    \item A visible or palpable midline bulge on straining, indicating rectus diastasis
    \item Persistent fullness of the lower abdomen that resists weight loss
    \item Skin irritation, intertrigo, or rashes within the abdominal skin folds
    \item A stable body weight, because results are eroded by further weight change or pregnancy
\end{itemize}

\section*{Differential Diagnosis}
Pre-operative assessment distinguishes conditions that change the surgical approach.

\begin{itemize}
    \item Excess skin requiring excision versus predominantly fatty tissue treatable by liposuction alone
    \item True rectus diastasis versus normal muscle tone and subcutaneous fat
    \item An abdominal wall hernia versus simple stretching of the musculature
    \item Weight that is stable and appropriate for surgery versus the need for weight loss first
    \item Realistic expectations achievable by the procedure versus goals the operation cannot meet
\end{itemize}

\section*{When to Seek Medical Care}
Before surgery, a full medical assessment identifies conditions such as hypertension, diabetes, or an elevated risk of blood clots that should be optimised first. After surgery, the surgical team should be contacted urgently if any of the following develop.

\begin{itemize}
    \item Fever that does not settle
    \item Spreading redness, warmth, discharge, or an opening wound
    \item Sudden or heavy bleeding from the wound
    \item Severe pain not controlled by prescribed analgesia
    \item Pain and swelling in one calf, chest pain, or difficulty breathing, which may indicate a blood clot
    \item Wound separation or skin that becomes dark and dusky, suggesting compromised blood supply
\end{itemize}

\textbf{Contact your local emergency services for:} collapse, severe bleeding, or sudden severe chest pain or breathlessness.

\section*{Diagnosis and Investigations}
Assessment is primarily clinical, with investigations directed at surgical fitness.

\begin{itemize}
    \item \textbf{History:} weight stability, pregnancy history, previous abdominal surgery, smoking, and current medications
    \item \textbf{Examination:} assessment of skin quality, fat distribution, and the abdominal wall, including a check for hernias and the degree of rectus separation
    \item \textbf{Body mass index:} surgery is safest at a stable, moderate weight
    \item \textbf{Photography:} standardised images for surgical planning and for comparison after recovery
    \item \textbf{Pre-operative tests:} blood tests and medical review tailored to the individual's risk
    \item \textbf{Expectation discussion and informed consent:} covering the scar, recovery time, and realistic outcomes
\end{itemize}

\section*{Management}
The specific technique is selected according to the amount of skin and fat and the location of the excess.

\begin{itemize}
    \item \textbf{Full abdominoplasty:} a low transverse incision, elevation of the skin and fat flap, tightening of the rectus muscles, excision of excess skin, and repositioning of the umbilicus
    \item \textbf{Mini abdominoplasty:} a shorter procedure for excess skin confined largely to the area below the umbilicus
    \item \textbf{Combined procedures:} liposuction of the flanks may be added, or abdominoplasty combined with other body-contouring surgery after major weight loss
    \item \textbf{Anaesthesia:} generally a general anaesthetic
    \item \textbf{Drains and compression garments:} used to reduce fluid collections and postoperative swelling
    \item \textbf{Pain relief and wound care} in the days following surgery
    \item \textbf{Gradual resumption of activity}, guided by the surgical team
    \item \textbf{Scar management:} silicone sheets or gels and sun protection while the scar matures
\end{itemize}

\section*{Complications}
\begin{itemize}
    \item Seroma, a collection of fluid beneath the skin, which is the most common complication and may require drainage
    \item Haematoma, or accumulation of blood within the wound
    \item Wound infection
    \item Poor wound healing, skin necrosis, or separation of the wound edges
    \item Permanent numbness or altered sensation across the lower abdomen
    \item Weakness of the abdominal wall or a new incisional hernia
    \item Venous thromboembolism involving the legs or lungs, a recognised risk of trunk surgery
    \item A long, permanent scar and possible contour irregularities
    \item The need for a revision procedure to address unsatisfactory results
\end{itemize}

\section*{Prognosis and Outcomes}
Most people are satisfied with both the appearance and the functional benefits of abdominoplasty, and the results remain stable provided weight is maintained. Typical recovery takes several weeks, with swelling settling over months and the scar fading gradually over a year or more. Significant weight gain or a subsequent pregnancy can reduce or reverse the improvements, which is why surgery is best planned once weight is stable and no further pregnancies are intended.

\section*{Prevention and Self-Care}
\begin{itemize}
    \item Reach and maintain a stable, healthy weight before surgery
    \item Do not smoke, as smoking greatly increases the risks of wound complications and blood clots
    \item Choose an experienced, appropriately qualified specialist surgeon
    \item Hold realistic expectations about what the procedure can achieve
    \item Follow the surgical team's advice on drains, dressings, and activity levels
    \item Wear compression garments and take prescribed medicines as directed
    \item Protect the scar from sun exposure and attend all follow-up appointments
    \item Avoid pregnancy or substantial weight gain in the period after surgery
\end{itemize}

\section*{Sources and Further Reading}
The following organisations provide authoritative information for patients considering abdominoplasty and for students learning about the procedure.

\begin{itemize}
    \item Mayo Clinic. \textit{Tummy tuck (abdominoplasty): overview and what to expect.} https://www.mayoclinic.org/tests-procedures/tummy-tuck/
    \item NHS. \textit{Cosmetic procedures: tummy tuck (abdominoplasty).} https://www.nhs.uk/conditions/cosmetic-procedures/
    \item MSD Manual. \textit{Plastic and reconstructive surgery overview.} https://www.msdmanuals.com/
    \item MedlinePlus. \textit{Abdominoplasty (tummy tuck) health information.} https://medlineplus.gov/ency/article/002981.htm
    \item American Society of Plastic Surgeons. \textit{Tummy tuck procedure information.} https://www.plasticsurgery.org/
    \item Royal Australasian College of Surgeons. \textit{Patient information on cosmetic surgery.} https://www.surgeons.org/
\end{itemize}
